\documentclass{article}

% CAISc 2026 official style (preprint = non-anonymous, for arXiv/sharing)
\usepackage[preprint]{caisc_2026}

% Packages
\usepackage[utf8]{inputenc}
\usepackage{amsmath,amssymb,amsfonts}
\usepackage{graphicx}
\usepackage{booktabs}
\usepackage{array}
\usepackage{hyperref}
\usepackage{xcolor}

\hypersetup{
  colorlinks=true,
  linkcolor=black,
  citecolor=black,
  urlcolor=black
}

% Spacing tweaks
\emergencystretch=2em
\setlength{\floatsep}{8pt plus 2pt minus 2pt}
\setlength{\textfloatsep}{10pt plus 3pt minus 3pt}
\setlength{\intextsep}{8pt plus 2pt minus 2pt}

% Custom commands
\newcommand{\sgmvpi}{$(S, G, E, V, M, \pi)$}
\newcommand{\framework}{\textsc{Servo}}

% Title
\title{\textbf{Formalizing AI Scientist Systems:\\
A Unified Framework for Automated Scientific Discovery}}

\author{
  Jaesol Shin\thanks{This manuscript was drafted with substantial AI assistance (Claude Code, Anthropic). The human author provided conceptual direction, framework design, and final judgment. See the AI Involvement Checklist for full details.}\\
  \small{WeDataLab}\\
  \texttt{ysys143@wedatalab.com}
}

\date{May 2026}

\begin{document}

\maketitle

\begin{abstract}
The emergence of end-to-end systems capable of automating hypothesis generation, experiment execution, and knowledge synthesis represents a significant change in how science is conducted. Despite rapid empirical progress, \emph{AI Scientist} systems lack a unifying formal framework, which makes principled comparison across architectures and scientific domains difficult. We propose \framework{}, a six-component framework \sgmvpi{} comprising the \textbf{S}earch space, hypothesis \textbf{G}enerator, \textbf{E}xperiment executor, \textbf{V}alidator, \textbf{M}emory, and search \textbf{P}olicy ($\pi$), grounded in the theory of Partially Observable Markov Decision Processes (POMDPs) and Bayesian Experimental Design (BED). We apply this framework to analyze core AI Scientist systems and their instantiations across seven scientific domains (mathematics, physics, chemistry, materials science, biology, social science, and computer science), and separately examine a complementary artifact infrastructure proposal that addresses the output format and verification layer of such systems. Our analysis surfaces three persistent open problems: (1) the non-automatability of novelty evaluation, which remains dependent on human judgment despite multi-layer automated proxies; (2) the absence of BED-optimal search policies, despite the theoretical availability of tractable approximations; and (3) the underdevelopment of semantic memory $M_s$, whose absence prevents accumulated experimental experience from improving future search. We demonstrate that the completeness of the validator $V$ is the single most consequential design variable: it determines whether closed-loop autonomous operation is feasible in each domain, and current AI Scientist systems can be read as a progression toward more complete validators.
\end{abstract}

\section{Introduction}

AI Scientist systems---agents that independently generate hypotheses, execute experiments, and synthesize knowledge---have proliferated rapidly~\citep{boiko2023emergent,lu2024aiscientist,lu2026aiscientist,schmidgall2025agentlab,google2026aletheia,liu2026lasthuman}, yet the field lacks a shared formal vocabulary. Systems range from tool-augmented synthesis agents (Coscientist~\citep{boiko2023emergent}) through closed-loop manuscript-writing pipelines (The AI Scientist~\citep{lu2024aiscientist,lu2026aiscientist}) to formal-mathematics agents (Aletheia~\citep{google2026aletheia}), each introducing bespoke terminology that makes systematic cross-system comparison difficult. Without a shared decomposition, questions that bear directly on deployment remain unanswered: \textit{What determines whether a system can operate in a closed loop? Why does autonomous novelty evaluation fail consistently? What is the theoretically optimal search policy?}

We propose \framework{} \sgmvpi{}, a formal framework grounding AI Scientist systems in POMDP and Bayesian Experimental Design theory (Section~\ref{sec:framework}). We apply it to analyze four core systems, an artifact infrastructure proposal, and instantiations across seven scientific fields (Sections~\ref{sec:core},~\ref{sec:domains}); and identify three open problems with falsification criteria (Section~\ref{sec:openproblems}). The analysis reveals a recurring pattern: validator ($V$) completeness determines closed-loop feasibility, not the sophistication of $G$ or $\pi$.

\noindent\textbf{Scope.} We cover systems whose primary goal is automated scientific \emph{discovery} through hypothesis-experiment-validation cycles, excluding specialized tools within human-conducted pipelines~\citep{jumper2021alphafold}.

\section{Background}
\label{sec:background}

Following~\citet{kaelbling1998pomdp}, we model AI Scientist systems as \emph{partially observable Markov decision processes}: the true state $s$ (unknown structure of nature) is inaccessible; the system maintains a belief $b(s)$ updated by experimental observations. Formally, an AI Scientist POMDP is the tuple $(\mathcal{S}, \mathcal{A}, T, R, \Omega, \mathcal{O})$, where $\mathcal{A}$ is the action space (hypothesis generation and experiment execution), $T$ is the transition function, $R$ is the reward (scientific value of discovery), and $\mathcal{O}$ is the noisy observation process. Three properties distinguish this from standard control: $R$ cannot be pre-specified (significance is not defined a priori); $\mathcal{S}$ is unbounded; and $T$ is non-stationary (accumulated knowledge changes the exploration frontier). Bayesian Experimental Design~\citep{chaloner1995bayesian,rainforth2024modern} provides the normative criterion for the search policy $\pi$: select $d^*$ maximizing Expected Information Gain (EIG),

\begin{equation}
  d^* = \arg\max_{d \in \mathcal{D}}\; \mathrm{EIG}(d) = \arg\max_{d}\; \mathbb{E}\left[H(\theta) - H(\theta \mid y, d)\right]
  \label{eq:eid}
\end{equation}

where $H(\theta)$ is prior entropy over parameters $\theta$ and $H(\theta \mid y, d)$ is expected posterior entropy after observing $y$ under $d$. EIG can be approximated via variational methods~\citep{foster2019variational}, amortized policies~\citep{foster2021dad}, and RL~\citep{blau2022rl}, yet no current AI Scientist system implements it---a gap we identify in Section~\ref{sec:openproblems}.

\section{Related Work}
\label{sec:related}

\textbf{AI Scientist systems.} Prior work~\citep{lu2024aiscientist,lu2026aiscientist,schmidgall2025agentlab,boiko2023emergent} provides qualitative characterizations of individual systems but no shared framework for cross-system comparison; idea-generation studies~\citep{si2024novelideas,baek2024researchagent} establish quality baselines without unifying the full discovery loop.

\textbf{Theoretical foundations.} The POMDP framework~\citep{kaelbling1998pomdp} and BED~\citep{chaloner1995bayesian,rainforth2024modern} are classical; we are the first to explicitly apply BED as a design target for AI Scientist search policies, connecting POMDP (structural formulation) to BED (normative policy criterion). The reward specification and specification gaming literatures study the consequences of unreliable reward signals; \framework{} extends this lens to AI Scientist systems by stratifying $V$ into reliability levels and connecting each level to specific closed-loop failure modes.

\textbf{Formal models and artifact infrastructure.} \citet{taniguchi2024cpc} propose Collective Predictive Coding (CPC-MS) as a social model of science, recasting scientific activity as decentralized Bayesian inference across a community. CPC-MS and \framework{} are complementary: CPC-MS provides a \emph{collective} and \emph{social} model in which peer review is a Metropolis-Hastings consensus process; \framework{} provides the \emph{single-system} and \emph{modular} decomposition that must be instantiated at each node. The three open problems we identify are precisely the necessary conditions for a CPC-MS-compliant collective system to be realizable. \citet{liu2026lasthuman} address the artifact layer --- output format and verification protocol --- rather than the discovery loop itself; we position ARA as a $V$ and $M_s$ infrastructure contribution within \framework{}, discussed in Sections~\ref{sec:ara} and~\ref{sec:openproblems}.

\section{The \framework{} Framework}
\label{sec:framework}

We define an AI Scientist system as a tuple \sgmvpi{} with the following components:

\begin{equation}
  \text{\framework{}} = (S,\; G,\; E,\; V,\; M,\; \pi)
\end{equation}

\begin{description}
  \item[$S$: Search Space] The set of candidate hypotheses, programs, molecules, experimental conditions, or other objects subject to exploration. $S$ may be static or dynamically expanding (open-ended).

  \item[$G: M \times S \rightarrow h$: Hypothesis Generator] A function that proposes new candidate $h \in S$ conditioned on accumulated memory $M$ and the current state of the search space. $G$ may involve LLM generation, evolutionary mutation, or symbolic search.

  \item[$E: h \times P \rightarrow o$: Experiment Executor] A function that executes hypothesis $h$ under protocol $P$ to produce observation $o$. The fidelity of $E$ ranges from fast computational simulation (low cost, potential surrogate error) to physical wet-lab execution (high cost, ground truth).

  \item[$V: o \rightarrow r$: Validator] A function that assigns a reward signal $r$ to observation $o$. We distinguish four layers of $V$:
    \begin{align*}
      V_\text{syntax}   &: o \rightarrow \{\text{pass}, \text{fail}\} & \text{(deterministic, fast)}\\
      V_\text{semantic} &: o \rightarrow r \in [0,1]                  & \text{(LLM-based, stochastic)}\\
      V_\text{empirical}&: o \rightarrow r                            & \text{(reproduction-based)}\\
      V_\text{human}    &: o \rightarrow \{\text{novel}, \text{significant}\} & \text{(non-automatizable)}
    \end{align*}

  \item[$M$: Memory] Structured as $(M_e, M_s, M_p)$. $M_e$ (episodic) stores individual records $(h_i, o_i, r_i)$; $M_s$ (semantic) stores distilled knowledge patterns; $M_p$ (procedural) stores learned protocols. The structure of $M$ directly determines the quality of $G$ and $\pi$.

  \item[$\pi: b(S) \times M \rightarrow h$: Search Policy] The decision rule for selecting the next hypothesis. The belief state $b(S)$ is updated after each experiment. In the BED interpretation, $\pi^* = \arg\max_d \mathrm{EIG}(d)$ (Equation~\ref{eq:eid}). In practice, current systems use policies ranging from reactive (ReAct) to greedy to tree search.
\end{description}

\noindent The operational loop is:
\begin{align}
  \pi(b, M) &\;\rightarrow\; G(M_s, S) \;\rightarrow\; E(h, P) \;\rightarrow\; V(o) \nonumber\\
            &\;\rightarrow\; M \mathrel{+}= (h, o, r) \;\rightarrow\; b(S) \text{ update} \;\rightarrow\; \text{repeat}
\end{align}

We further introduce $H \in [0,1]$ (distinct from the entropy function $H(\theta)$ in Eq.~\ref{eq:eid}) as an auxiliary parameter denoting the degree of human intervention: $H=0$ for fully autonomous systems and $H=1$ for systems in which humans assume the role of $G$ or $\pi$.

\noindent\textbf{V-completeness as a design dimension.} The four $V$ layers form a partial order under reliability and automatability. $V_\text{syntax}$ is fully automatable and deterministic; $V_\text{semantic}$ is automatable but stochastic and subject to LLM biases; $V_\text{empirical}$ varies from well-calibrated (DFT for crystal stability) to poorly calibrated (GNN surrogates for wet-lab outcomes); $V_\text{human}$ is irreplaceable for novelty and significance but is the rate-limiting step for loop frequency. In domains where correctness is mechanically decidable (formal proof kernels, type-checkers), $V_\text{formal}$ --- a fifth special case --- achieves the ideal limit of $V_\text{empirical}$: perfect calibration and zero false positives. Table~\ref{tab:domain-comparison} uses $V_\text{formal}$ for mathematics (formal) to mark this distinct reliability regime. We define \emph{$V$ completeness} as the highest layer of $V$ that is both accessible and reliably calibrated in a given system and domain. Accessibility without calibration does not count: Agent Laboratory's $V_s$ (LLM reviewer) is accessible but systematically biased, so it contributes imperfect $V$ completeness rather than full $V_s$ completeness. This dimension is the primary design variable because it sets a ceiling on what $\pi$ can learn and what $G$ can be rewarded for: without a reliable $V$, no policy can distinguish progress from regression, and no generator can improve with experience. The cross-domain pattern in Section~\ref{sec:domains} is a direct consequence: every domain's loop status is determined by $V$ completeness, not by the sophistication of its $G$ or $\pi$.

\section{Analysis of Core AI Scientist Systems}
\label{sec:core}

We apply \framework{} to four end-to-end systems~\citep{boiko2023emergent,lu2024aiscientist,lu2026aiscientist,schmidgall2025agentlab} and one artifact infrastructure proposal. Table~\ref{tab:core-comparison} synthesizes the key dimensions; the cross-system pattern is that every architectural advance is a gain in $V$ completeness, and $V$ completeness is what unlocks closed-loop operation.

\vspace{6pt}
\begin{table}[h]
\centering
\small
\setlength{\tabcolsep}{4pt}
\begin{tabular}{@{}lcccc@{}}
\toprule
& \textbf{Coscientist} & \textbf{AI Sci.\ (2024)} & \textbf{AI Sci.\ (2026)} & \textbf{AgentLab} \\
\midrule
Novelty measure   & None         & LLM self-score & 3-layer         & Human-delegated \\
Loop closed?      & No           & Yes (greedy)   & Yes (tree)      & Partial         \\
$\pi$ type        & ReAct        & Greedy         & Tree search     & Human/seq.      \\
$V$ completeness  & $V_h$ only   & $V_s$          & $V_e+V_s+V_h$  & $V_s$ (biased)  \\
$M$ structure     & Ephemeral    & $M_e$          & $M_e+M_s$       & $M_e$           \\
$H$ (approx.)     & G (${\approx}1$) & Low (${\approx}0$) & Low (${\approx}0$) & G+$\pi$ (${\approx}0.8$) \\
\bottomrule
\end{tabular}
\caption{Comparative \framework{} analysis of four core AI Scientist systems. $V_s$=semantic, $V_e$=empirical, $V_h$=human.}
\label{tab:core-comparison}
\end{table}

\noindent The progression across the four systems is a systematic increase in $V$ completeness, not in $G$ or $\pi$ sophistication. Coscientist's single GC-MS measurement ($V_h$ only) prevents loop closure entirely. The AI Scientist (Sakana) adds $V_s$ via a 5-ensemble GPT-4o reviewer (balanced accuracy 0.65 vs.\ human 0.66~\citep{lu2024aiscientist}), enabling the first genuinely closed loop despite greedy $\pi$. The AI Scientist (Nature) stacks $V_e$ (replication nodes with mean $\pm$ s.d.), $V_s$ (automated reviewer, balanced accuracy 0.69~\citep{lu2026aiscientist}), and actual peer review ($V_h$ via ICLR workshop), achieving the most complete $V$ in the class. Agent Laboratory delegates $V$ to $V_s$ (NeurIPS-guideline LLM reviewer) but documents $+$2.3-point systematic over-estimation relative to human PhD students---illustrating that adding a $V$ layer does not guarantee $V$ completeness if the layer is unreliable. The three systems that close the loop each do so via a $V$ upgrade; Agent Laboratory, which adds a $V_s$ layer without calibrating it, demonstrates the converse --- an accessible but unreliable $V$ does not close the loop. No system in this class achieves improved loop feasibility through $G$ or $\pi$ advances alone, without a corresponding $V$ improvement.

\subsection{Artifact Infrastructure: Agent-Native Research Artifacts}
\label{sec:ara}

\citet{liu2026lasthuman} address a complementary problem: what the \emph{output} of the AI Scientist loop should look like and how it can be independently verified. Agent-Native Research Artifacts (ARA) replace the narrative paper with a four-layer ontology: \texttt{/logic} (claims with falsification criteria), \texttt{/src} (executable code), \texttt{/trace} (exploration graph: question, decision, experiment, dead\_end, pivot), and \texttt{/evidence} (raw results). The ARA Seal provides three-level automated verification---$V_\text{syntax}$ (schema conformance), $V_\text{semantic}$ (Rigor Auditor, 82.6\% average detection rate across five structural violation types), and $V_\text{empirical}$ (sandboxed reproduction)---with $V_\text{human}$ reserved for significance, novelty, and taste: \textit{``judgments of significance, novelty, and taste are reserved for human reviewers''}~\citep{liu2026lasthuman}. The motivation is quantitative: only 45.4\% of reproduction requirements are fully specified in published papers, and 90.2\% of extension costs are consumed by failed exploration~\citep{liu2026lasthuman}. The \texttt{/trace} layer addresses $M_e$ and $M_s$: preserved failure traces accelerated agent performance on extension tasks but also anchored more capable agents to prior solution paths---richer $M_e$ does not monotonically improve $\pi$.

\section{Domain Analysis}
\label{sec:domains}

Table~\ref{tab:domain-comparison} maps \framework{} instantiations across seven scientific domains; $V$ completeness co-varies with loop closability without exception. To illustrate scale: Aletheia~\citep{google2026aletheia} attempted 700 Erd\H{o}s problems; of the 200 solutions selected for human audit, 13 (6.5\%) were rated meaningfully correct; GNoME~\citep{merchant2023gnome} screened 2.2 million candidate crystal structures via high-throughput DFT. Systems analyzed span mathematics (formal:~\citep{romera2023funsearch,xin2024deepseek}; NL:~\citep{google2026aletheia}), physics~\citep{udrescu2020afeynman,cranmer2020symbolic}, chemistry~\citep{bran2023chemcrow,sprueill2024chemreasoner}, materials~\citep{merchant2023gnome}, biology~\citep{odonoghue2023bioplanner}, social science~\citep{manning2024automated,aher2023turing,park2023generative}, and CS/algorithms~\citep{real2020automlzero,sartori2024algodiscovery,jaber2026autokernel}.

\begin{table}[h]
\centering
\small
\setlength{\tabcolsep}{4pt}
\begin{tabular}{@{}lccccp{2.6cm}@{}}
\toprule
Domain & $V$ type & $V$ reliability & Loop & $\pi$ & Primary bottleneck \\
\midrule
Math (formal) & $V_\text{formal}$ & Perfect oracle & Fully closed & RL/MCTS  & Pretraining contamination \\
Math (NL) & $V_\text{sem}+V_h$ & Imperfect & Closed (low quality) & Seq.\ refine & Hallucination; plagiarism \\
Physics & $V_\text{emp}+V_h$ & Medium & None/partial & Deterministic & No cross-problem $M$ \\
Chemistry & $V_\text{emp}$ (surr.) & Medium & Comp.\ only & Beam & Wet-lab loop open \\
Materials & $V_\text{emp}$ (DFT) & High & Comp.\ closed & Uncert.\ samp. & DFT wall-time; synth.\ lag \\
Biology & $V_\text{emp}+V_h$ & Low--med & Partial & None & $V$ metrics insufficient \\
Social Sci. & Stat.\ test & Low--med & Partial & Factorial & Circular validity \\[3pt]
CS / Algo. & $V_\text{emp}$ & High & Fully closed & Evol./LLM & Interpretability absent \\
\bottomrule
\end{tabular}
\caption{Cross-domain \framework{} analysis across seven domains. $V$ reliability refers to agreement with ground truth.}
\label{tab:domain-comparison}
\end{table}

\noindent The domain pattern is not coincidental --- it is structurally determined by $V$ type. \textbf{Mathematics (formal):} $V_\text{formal}$ is binary and unambiguous; RL is stable because reward hacking is impossible by construction, and both FunSearch and DeepSeek-Prover achieve fully closed loops. \textbf{Mathematics (NL):} switching to natural-language proof spaces immediately reintroduces $V_\text{human}$ as the only reliable oracle; Aletheia's 6.5\% meaningful-correctness rate among audited solutions confirms the practical gap. \textbf{Physics:} symbolic interpretability (``this formula is new to cosmologists'') is a judgment that no current $V$ can substitute; each problem restarts from scratch because there is no cross-problem $M_s$. \textbf{Chemistry and biology:} the computational loop closes (GNN surrogate or protocol-accuracy $V$) but the physical loop does not; NMR, MS, and wet-lab analytical characterization remain human-delegated, and the wet-lab loop is structurally open at the measurement boundary. \textbf{Materials science:} GNoME's DFT oracle is the field's closest approximation to an ideal $V$---well-calibrated, tractable, and noise-free within its domain; it is not coincidental that this is the single domain where $\pi$ approximates BED theory. \textbf{Social science:} when $G$ and $E$ share the same underlying model, the evaluator systematically confirms the generator's implicit biases~\citep{aher2023turing}; hyper-accuracy distortion means LLM-simulated subjects are too accurate to simulate real human population distributions. \textbf{CS/algorithms:} $V_\text{emp}$ (runtime speedup, validation accuracy) is high-reliability and enables fully closed loops, but interpretability of discovered algorithms remains absent --- knowing \emph{that} an algorithm works is distinct from understanding \emph{why}.


\section{Open Problems}
\label{sec:openproblems}

Our analysis identifies three open problems that persist across all systems and domains.

\subsection{The Non-Automatability of Novelty}

Novelty cannot be reliably evaluated by automated means: every system from Coscientist (no mechanism) to The AI Scientist Nature (three-layer validation) converges on this conclusion. Two failure modes are documented: LLM self-evaluation bias (systematic over-estimation; hallucinated results) and pretraining contamination (systems may ``discover'' results implicit in training data, as Aletheia documents for Erd\H{o}s problems). The ARA Seal~\citep{liu2026lasthuman} demarcates the automated frontier: its Rigor Auditor achieves 82.6\% average detection across five structural violation types, but detects only 22\% of orphaned-experiment violations specifically (the two rates measure different denominators and are not complements). It cannot cross the novelty boundary. Concretely: Agent Laboratory's LLM reviewer over-estimated paper quality by $+$2.3 points on the NeurIPS review scale relative to human PhD students; The AI Scientist (Sakana) hallucinated ablation tables; and of the 200 Erd\H{o}s solutions selected for human audit (from 700 attempted), only 13 (6.5\%) were rated meaningfully correct despite surfacing many superficially plausible candidates.

\textit{Falsification criterion:} Refuted if an automated $V$ achieves human-level novelty discrimination (inter-rater agreement $\geq$0.80) on a contamination-controlled benchmark spanning $\geq$3 distinct scientific domains.

The deeper obstacle is that novelty, significance, and correctness are distinct properties that current systems treat as a single scalar. The AI Scientist~\citep{lu2024aiscientist} illustrates this directly: its 1--10 ``novelty'' score collapses technical correctness (does the code run?), archive novelty (does it differ from past work?), and scientific significance (is it worth pursuing?) into one number. Each property has a different evaluator, a different reference corpus, and a different contamination sensitivity; a system that maximizes the composite score may satisfy none of the three individually. A viable automated $V$ for novelty requires at minimum: (a) a contamination-controlled reference corpus that separates what the model was trained on from what the field has published; (b) architectural separation between $G$ and $V$ so the evaluating model has not seen its own outputs; and (c) a formal definition of ``novel to the field'' that does not reduce to ``not lexically in training data.'' None of these conditions hold in any current AI Scientist system.

\subsection{The BED--Practice Gap for $\pi$}

Equation~\ref{eq:eid} specifies the theoretically optimal policy, yet no current system implements it: all deployed policies are heuristic (ReAct, greedy, tree search, human-guided). The partial exception is GNoME~\citep{merchant2023gnome}, whose deep-ensemble uncertainty sampling approximates EIG over a tractable GNN posterior---but only because DFT provides a well-calibrated $V$ and the hypothesis space admits a learnable probabilistic model, conditions absent from LLM-based systems where the ``posterior'' is implicit in model weights~\citep{foster2021dad}. Concretely, every system in Table~\ref{tab:core-comparison} degrades to heuristic $\pi$: Coscientist selects the next tool reactively, The AI Scientist (Sakana) greedily avoids archive similarity, and Agent Laboratory follows a fixed sequential schedule---none improve with additional experimental data, the failure mode EIG-optimal $\pi$ would avoid.

\textit{Falsification criterion:} Refuted if any AI Scientist system implements a policy that explicitly computes and maximizes EIG over a posterior $p(\theta \mid y)$ in a domain where the posterior is not tractable from first principles---ruling out tractable-model proxies such as GNoME's deep-ensemble policy.

The tractability gap has a precise architectural characterization. EIG requires computing $\mathbb{E}_y[H(\theta) - H(\theta \mid y, d)]$, which in turn requires a posterior $p(\theta \mid y)$ over hypothesis parameters $\theta$ given experimental outcomes $y$. In LLM-based systems, $\theta$ is implicit in frozen model weights; there is no mechanism to update these weights per experiment, and no tractable approximation to the posterior entropy over the implied hypothesis space. Two paths exist. One is to maintain an explicit probabilistic surrogate of the hypothesis space---GNoME's approach---grounding EIG in the surrogate's posterior rather than model weights. This works where a learnable, calibrated input-outcome mapping exists, which excludes most LLM-based domains. The other is amortized EIG estimation~\citep{foster2021dad}, bypassing explicit posteriors; this remains unsolved for open-ended domains where $\mathcal{S}$ is not fixed in advance.

\subsection{The Underdevelopment of Semantic Memory $M_s$}

Current systems store raw experiment records ($M_e$) but do not distill them into transferable patterns ($M_s$): a system with 1,000 experiments does not search more intelligently than one with 10. The ARA \texttt{/trace} layer~\citep{liu2026lasthuman} is the closest approximation, but traces accelerated extension performance while also anchoring capable agents to prior solution paths---$M_s$ structure must be paired with selective forgetting mechanisms to be beneficial. Concretely, GNoME's GNN---one of the few systems where $G(M_s, S)$ genuinely improves---accumulates domain-specific thermodynamic knowledge after 460K DFT calculations, yet this knowledge does not transfer to chemistry or biology; the system is no more capable in any domain it has not explicitly trained on.

\textit{Falsification criterion:} Refuted if accumulated cross-experiment knowledge demonstrably improves $G$-generated hypothesis quality across task instances in a domain where the system was not explicitly pre-trained or purpose-built, measured by a pre-specified metric.

The set-union coverage model---where memory sharing is perfect and additive---is the trivially ideal limit of $M_s$: knowledge simply accumulates, and the only design question is connectivity. The genuinely hard $M_s$ problems arise when memory is \emph{lossy} (semantic retrieval degrades with distance from training distribution), \emph{compressed} (distillation loses evidence the original agent deemed unimportant), \emph{conflicting} (two agents may reach contradictory conclusions from overlapping experiments), or \emph{strategically filtered} (incentives to withhold negative results). In each case, accumulated experience can degrade $G$'s performance: a system trained on its own noisy memory may converge prematurely, systematically over-represent successful strategies, or amplify false positives from early exploration. GNoME demonstrates that within-domain $M_s$ accumulation \emph{can} work: the GNN genuinely improves as DFT outcomes accumulate. The open problem begins where GNoME ends---the same memory structure cannot be reused in chemistry or biology without full retraining, and even within materials science it degrades under distribution shift. Building $M_s$ mechanisms that generalize across experimental contexts, forget selectively, and remain robust to conflicting evidence is a necessary precondition for AI Scientist systems whose $G(M_s, S)$ improves with scale rather than plateauing.

\noindent The three open problems are not independent. A system with complete $V$ and tractable $M_s$ would have the necessary prerequisites for an EIG-optimal $\pi$: complete $V$ provides the calibrated reward signal that EIG requires, and tractable $M_s$ provides the hypothesis-space model over which EIG is computed. Whether these prerequisites are sufficient for tractable EIG estimation in non-stationary or unbounded $S$ remains open, but they are necessary. They therefore converge on a single structural gap, and progress on any one unlocks disproportionate gains toward the other two. They are also precisely the necessary conditions for a collective AI co-scientist~\citep{taniguchi2024cpc}: $V_\text{collective}$ requires solved novelty automatability; a principled collective $\pi$ requires closed BED--practice gap; and the shared semantic memory $w$ requires developed $M_s$. ARA~\citep{liu2026lasthuman} may be the most tractable entry point, simultaneously externalizing $M_s$ and calibrating the sub-novelty layers of $V$.

\section{Discussion}
\label{sec:discussion}

\textbf{Design hierarchy.} Our analysis establishes a clear priority ordering: $V$ design precedes all other decisions, because it determines whether closed-loop operation is feasible and what form $\pi$ can take. In domains where $V$ is noisy or fallible (chemistry, biology, social science), even EIG-optimal $\pi$ and perfectly structured $M_s$ cannot close the loop---the reward signal is too unreliable to sustain RL or guide principled search. Conversely, in mathematics (formal), $V_\text{formal}$ is so well-specified that even greedy $G$ and simple $\pi$ achieve closed loops. The practical implication: before optimizing $G$ or $\pi$, designers should ask what level of $V$ completeness the domain actually supports.

\textbf{Human oversight.} Rather than treating $H \in [0,1]$ as a limitation to be minimized, our analysis suggests that explicit human integration at the $V_\text{human}$ and $\pi$ levels is both practically necessary and theoretically appropriate. Human judgment provides access to an information source that no automated proxy currently replicates---particularly for the non-automatizable components of novelty and significance. Systems that design $H$ explicitly (as ARA does by reserving significance and novelty for human reviewers) may achieve better long-run outcomes than those that attempt to minimize it. The goal is not full autonomy but \emph{appropriate} autonomy: closing the loop at the layers where $V$ is reliable while preserving human judgment at the layers where it is not.

\textbf{Towards collective AI co-scientists.} All systems analyzed operate as single-agent loops. A natural extension is a collective AI scientist: a multi-agent system $\{\framework_k\}_{k=1}^K$ sharing an external semantic memory $w$ and coordinating through a collective validator $V_\text{collective}$~\citep{taniguchi2024cpc}. The three open problems we identify are precisely the necessary conditions for such a system: $V_\text{collective}$ requires solved novelty automatability; a principled collective $\pi$ requires closing the BED--practice gap; and the shared $w$ requires developed $M_s$. ARA~\citep{liu2026lasthuman} may be the most tractable entry point as a structured artifact protocol for $w$, simultaneously externalizing $M_s$ and calibrating the sub-novelty layers of $V$.

\textbf{Limitations.} \framework{} abstracts away several important factors. The cost structure of $E$ varies by orders of magnitude across domains (milliseconds for code execution, hours for DFT, weeks for wet-lab) and is not modeled. The non-stationarity of $S$ (open-ended problems) is treated as a property to be noted but not formalized. The single-agent framing abstracts away the social layer that \citet{taniguchi2024cpc} formalize; extending \framework{} to multi-agent settings is a natural direction for future work. Finally, the cross-domain comparison relies on publicly reported system descriptions; component characterizations may not capture all implementation details.

\section{Conclusion}
\label{sec:conclusion}

We introduced \framework{} \sgmvpi{}, a formal framework connecting AI Scientist systems to POMDP and BED theory, and applied it to four core systems, one artifact infrastructure proposal, and seven scientific domains. The central finding is that $V$ completeness is the primary determinant of closed-loop feasibility: every architectural advance across surveyed systems---from Coscientist's single-measurement $V_h$ to The AI Scientist Nature's three-layer validation---is ultimately a gain in $V$ completeness, not in $G$ sophistication or $\pi$ optimality.

The \framework{} decomposition provides a shared vocabulary that makes previously implicit cross-system comparisons explicit. Under this vocabulary: the reason mathematics (formal) has the most mature AI Scientist deployments is $V_\text{formal}$'s binary reliability, not domain simplicity; the reason chemistry and biology remain loop-open despite capable $G$ and $\pi$ implementations is $V$'s dependence on physical measurement; and the reason no system improves intelligently with scale \emph{beyond its specific training domain} is the absence of transferable $M_s$ mechanisms (GNoME's within-domain improvement is the exception that proves the rule: it works only because 460K DFT calculations were performed in a single, well-defined chemical space). These are structural constraints, not engineering gaps to be solved by better language models alone.

The three open problems---novelty non-automatability, the BED--practice gap, and $M_s$ underdevelopment---are not independent, and they are not equally tractable. ARA's Rigor Auditor represents meaningful progress on the sub-novelty layers of $V_\text{semantic}$; GNoME's uncertainty sampling demonstrates that EIG-adjacent $\pi$ is achievable when a tractable probabilistic model of the hypothesis space exists. These partial exceptions suggest the bottleneck is not fundamental impossibility but missing architectural components. What \framework{} offers is a vocabulary for naming them precisely: the field needs validators and memory structures that are both complete enough to close the loop and structured enough to support principled search.

\bibliographystyle{abbrvnat}
\bibliography{references}

% ---------------------------------------------------------------
\newpage

\section*{Conference For AI Scientists 2026 - AI Involvement Checklist}
\label{checklist_ai_involvement}

This checklist documents the role of AI systems at each stage of this research.
The checklist has two parts: a \textbf{Research Stage Assessment} (items 1--4) rating AI involvement and iteration effort, and an \textbf{AI System Documentation} section (items 5--6).

\subsection*{Research Stage Assessment}

\begin{enumerate}

  \item \textbf{Hypothesis development.}

    Answer: \involvementB{}

    Iteration Effort: \iterationMedium{}

    Explanation: The core \framework{} framework tuple \sgmvpi{}, the $V$-completeness thesis, the $H$ parameter, and the three open problems were conceived by the human author. A frontier LLM agent elaborated framework components, identified literature-level instantiations, and proposed theoretical connections. The human author directed all conceptual decisions and made several substantive corrections to the AI's elaborations.

  \item \textbf{Experimental design and implementation.}

    Answer: \involvementB{}

    Iteration Effort: \iterationMedium{}

    Explanation: This submission contains no computational experiments. The methodology for cross-system analysis (system selection criteria, component extraction protocol, domain comparison structure) was designed by the human author. A frontier LLM agent executed literature retrieval, extracted framework components from primary sources, and performed citation-level fact verification across four verification sweeps.

  \item \textbf{Analysis of data and interpretation of results.}

    Answer: \involvementC{}

    Iteration Effort: \iterationMedium{}

    Explanation: The LLM agent extracted \framework{} components for all analyzed systems from primary sources and performed citation-level fact verification (4 verification sweeps). The human author reviewed all extractions, corrected mischaracterizations, and synthesized the cross-domain pattern. The V-completeness thesis emerged from the human author's synthesis.

  \item \textbf{Writing.}

    Answer: \involvementC{}

    Iteration Effort: \iterationHigh{}

    Explanation: The LLM agent drafted the majority of the manuscript text through multiple revision cycles; the human author provided key conceptual passages (framework definitions in \S\ref{sec:framework}, core open problems in \S\ref{sec:openproblems}), directed structural decisions, reviewed all drafts, and made final editorial judgments.

\end{enumerate}

\subsection*{AI System Documentation}

\begin{enumerate}
  \setcounter{enumi}{4}

  \item \textbf{AI systems used.}

    Description: A frontier large language model with code execution and long-document processing capabilities was used as the primary research and writing assistant. The system operated as an interactive CLI agent with file-system access, able to read primary sources, execute and debug Python simulations, and iteratively revise LaTeX manuscripts.

  \item \textbf{Observed AI Limitations.}

    Description: (1)~Systematic citation fabrication: the agent generated plausible but non-existent quotes and statistics, requiring four dedicated citation-verification sweeps using parallel sub-agents. (2)~Analytical overconfidence: the agent initially overclaimed empirical significance for mechanically-entailed theoretical consequences; multiple adversarial review rounds were needed to correct framing and remove overclaims. (3)~Context drift: over long sessions, the agent occasionally regressed previously corrected errors, requiring explicit re-verification of integrity items before each commit.

\end{enumerate}

% ---------------------------------------------------------------
\newpage

\section*{Conference For AI Scientists 2026 - Reproducibility and Responsibility Checklist}
\label{checklist_reproducibility}

\begin{enumerate}

\item {\bf Claims}
    \item[] Question: Do the main claims made in the abstract and introduction accurately reflect the paper's contributions and scope?
    \item[] Answer: \answerYes{}
    \item[] Justification: The abstract and introduction explicitly state the three contributions (framework, cross-domain analysis, open problems). The V-completeness thesis is presented as a cross-domain empirical pattern supported by natural comparison, not a proven theorem, and the paper explicitly acknowledges this scope in \S\ref{sec:discussion}.

\item {\bf Limitations}
    \item[] Question: Does the paper discuss the limitations of the work performed by the authors?
    \item[] Answer: \answerYes{}
    \item[] Justification: Limitations are discussed in \S\ref{sec:discussion} (framework abstracts domain cost structures, social layer, non-stationarity) and throughout the domain analysis (\S\ref{sec:domains}).

\item {\bf Theory assumptions and proofs}
    \item[] Question: For each theoretical result, does the paper provide the full set of assumptions and a complete proof?
    \item[] Answer: \answerNA{}
    \item[] Justification: The paper makes no formal theorem claims. The V-completeness thesis is a cross-domain empirical pattern supported by natural comparison, not a proven theorem.

\item {\bf Experimental result reproducibility}
    \item[] Question: Does the paper fully disclose all information needed to reproduce the main experimental results?
    \item[] Answer: \answerYes{}
    \item[] Justification: This paper contains no computational experiments. The cross-domain analysis methodology (component extraction from primary sources) is fully described in \S\ref{sec:core} and \S\ref{sec:domains}; all analyzed papers are publicly accessible.

\item {\bf Open access to data and code}
    \item[] Question: Does the paper provide open access to data and code?
    \item[] Answer: \answerNo{}
    \item[] Justification: This submission contains no computational experiments. All analyzed papers are publicly accessible, and the framework can be applied following the component definitions in \S\ref{sec:framework}.

\item {\bf Experimental setting/details}
    \item[] Question: Does the paper specify all details necessary to understand the results?
    \item[] Answer: \answerYes{}
    \item[] Justification: The cross-domain analysis methodology is described in \S\ref{sec:core} and \S\ref{sec:domains}. No computational experiments were performed in this submission.

\item {\bf Experiment statistical significance}
    \item[] Question: Does the paper report error bars or statistical significance information?
    \item[] Answer: \answerNA{}
    \item[] Justification: This paper contains no computational experiments and reports no statistical results requiring error bars or significance tests.

\item {\bf Experiments compute resources}
    \item[] Question: Does the paper provide sufficient information on compute resources?
    \item[] Answer: \answerNA{}
    \item[] Justification: This submission contains no computational experiments. The meta-analysis requires only standard document processing and no specialized compute.

\item {\bf Research Ethics}
    \item[] Question: Does the research conform with the NeurIPS Code of Ethics, except where CAISc's AI authorship policies apply?
    \item[] Answer: \answerYes{}
    \item[] Justification: This is a meta-analysis of published systems with no human subjects, no proprietary data, and no safety risks. AI authorship is documented per CAISc policy.

\item {\bf Broader impacts}
    \item[] Question: Does the paper discuss both potential positive and negative societal impacts?
    \item[] Answer: \answerYes{}
    \item[] Justification: \S\ref{sec:discussion} addresses both the potential for accelerating scientific discovery and the risk of premature convergence under widespread AI Scientist deployment (particularly via the $M_s$ underdevelopment problem and novelty non-automatability).

\end{enumerate}

\end{document}
